\section{Quotient Groups and Rings}
\begin{theorem}
If $H$ is an invariant subgroup of $G$, the cosets of $H$ in $G$ form a group under the group product given by
\[ (gH) (g'H) = gg'H.\]
\end{theorem}
The product is well defined and unique, and is also a coset. \\

It is associative, since 
\[ ((gH)(g'H)) (g'' H) = gg'g''H = (gH)( (g'H)(g''H)).\]

The coset $H = (eH)$ is a unity since
\[ (gH) (eH) = geH =gH \mbox{   and  } (eH)(gH) = egH = gH.\]

The coset $gH$ has an inverse $g^{-1}H$, since
\[ (gH)(g^{-1}H) = gg^{-1}H = eH = H \mbox{   and   } (g^{-1}H)(gH) = H\] similarly.


The quotient group formed by the cosets of $H$ and $G$ is called the {\it Quotient Group of $H$ in $G$}, or sometimes the {\it Factor Group of $G$ relative to $H$}.

\section{Quotient Group Associated with a Homomorphism}

As we have already hinted, the subjects of invariant subgroups and homomorphisms of groups are closely connectedl there is in fact a 1-1 correspondence between subgroups of a given group $G$ and the essentially different 
homomorphism which have $G$ ax the object space, so that the study of either tops is similar to that of the other. \\

Suppose that $H$ is an invariant subgroup of $G$, so that the quotient group $G/H$ is defined. Consider the mapping $\theta: G \goes G/H$ defined by $g\theta = gH$, the coset containing $g$, which is of course an element of $G/H$.

\begin{theorem}
The mapping $\theta: G \goes G/H$ defined by $g\theta = gH$ is a homomorphism.
\end{theorem}

\begin{proof}
\begin{eqnarray*}
(gg')\theta &=& gg'H = gg' H H\cr
                &=& gH g'H \cr
                &=& (g\theta) (g'\theta)
\end{eqnarray*}
\end{proof}

\begin{corollary}
The kernel $\theta$ is $H$ and the image is $G/H$, so that $\theta$ is an epimorphism.
\end{corollary}

The homomorphism $\theta$ is called the Natural Homomorphism or the Canonical Homomorphism associated with $H$. It is essential that $H$ be an invariant subgroup, for otherwise $G/H$ is not defined and we cannot speak of the product of the cosets $(gH)(g'H)$.

Any homomorphism $\theta: G\goes G'$ is closely related to the natural homomorphsim associated with an invariant subgroup of $G$. 

\begin{theorem}
The kernel of $\theta: G \goes G'$ is an invariant subgroup of $G$. 
\end{theorem}

If $h$ is in the kernel $H$ and $g$ is any element of $G$, we have
\begin{eqnarray*}
(g^{-1} h g)\theta  &=& (g^{-1}\theta) (h\theta) (g\theta)\cr
&=& (g\theta)^{-1}) f (g\theta), \mbox{  where $f$ is the neutral element of $G'$}\cr
&=& (g\theta)^{-1}(g\theta) = f
\end{eqnarray*}

\begin{theorem}
The mapping defined by $gH \goes g\theta$ is a well-defined mapping of $G/H$ onto $G\theta$.
\end{theorem}
If $g_1$ and $g_2$ are in the same coset, then $g_1\theta = g_2\theta$. If $g_1$ and $g_2$ are in the same coset, then $g_2^{-1}g_1 \in H$. Hence
$(g_2\theta)^{-1}(g_1\theta) = f$, the neutral element in $G'$, since $H$ is in the kernel of $\theta$. Therefore $g_2\theta = g_1\theta$. 

Also  any element of $G\theta$ is the image of some $g\in G$ and so the mapping is onto $G\theta$, although not necessarily onto $G'$. 

\begin{theorem}
The mapping defined by $gH \goes g\theta$ is an isomorphism between $G/H$ and $G\theta$. 
\end{theorem}
We have that 
\begin{eqnarray*}
(g_1 H) (g_2 H) &=& g_1 g_2 H, \cr
(g_1\theta) (g_2\theta) &=& (g_1 g_2)\theta, 
\end{eqnarray*}
because $\theta$ is a homomorphism and $g_1 g_2 H \goes (g_1 g_2)\theta$. 

Since this mapping is onto $G\theta$, it is also an isomorphism provided it is 1-1, i.e., provided the kernel is the neutral element of $G/H$. Suppose
$gH \goes f$. Then $g\theta \goes f$ and so $g\in H$, so that $gH = H$, the neutral element of $G/H$. \\

Notice that the sets of elements of $G$ which map into the particular elements of $G\theta$ are the cosets of $G$ relative to the kernel $H$. \\

We have proved that to any invariant subgroup $H$ of $G$ there is a natural homomorphism from $G$ onto $G/H$, and conversely that any homomorphism of G is  isomorphic to
a natural homomorphism associated with its kernel. Therefore there is a 1-1 correspondence between the invariant subgroups and the essential different homomorphisms of $G$. By 
essentially different homomorphisms we mean those having different kernels: two homomorphisms having the same kernel are both equivalent to the same natural homomorphism and therefore exhibit the same structure.\\ 

The image space $G'$ is largely irrelevant : the subgroup $G\theta$ of $G'$ is isomorphic to $G/H$ but $G'$ may be larger. Also there is no reason why $G\theta$ should be invariant in $G'$. \\